\subsection{用辅助可逆机推出循环不等式}
\label{th:clausiusconstruction}
把一次循环与外界的热交换分成若干份，装置从温度 $T_k$ 的热库吸收有符号热量 $Q_k$，放热时 $Q_k<0$。另选温度为 $T_0$ 的参考热库。对每个 $T_k$ 配一台辅助可逆机，使它从该热库取得的热量恰为 $-Q_k$；于是原装置与辅助机在 $T_k$ 热库的热交换完全抵消。

设该辅助机从参考热库吸热为 $Q_{0k}$。辅助机可正向也可反向运行，其可逆热量关系用有符号形式写为
\[
 -\frac{Q_k}{T_k}+\frac{Q_{0k}}{T_0}=0,
 \qquad Q_{0k}=T_0\frac{Q_k}{T_k}.
\]
全部工质均作循环，组合装置只与参考热库交换净热
$Q_0=T_0\sum_kQ_k/T_k$。第一定律给组合输出功 $W_{\rm comb}=Q_0$。若 $\sum_kQ_k/T_k>0$，组合的唯一效果便是从单一热库吸热并全部转成功，违反开尔文表述。因此
\[
 \sum_k\frac{Q_k}{T_k}\le0.
\]
把热交换分得足够细便得到循环积分形式。若原循环也可逆，整套装置可以反向运行，同样的不等式施于反循环，便与原式合成等号。

\Diagram{fig:auxiliaryclausius}{辅助可逆机抵消各分温热库的交换}{
 \node[draw](t)at(0,2){热库 $T_k$};
 \node[draw](a)at(-1.6,0){原循环};
 \node[draw](b)at(2,0){辅助可逆机};
 \node[draw](z)at(2,-1.8){参考热库 $T_0$};
 \draw[axis](t)--(a)node[midway,left]{$Q_k$};
 \draw[axis](b)--(t)node[midway,right]{$Q_k$};
 \draw[axis](z)--(b)node[midway,right]{$Q_{0k}$};
 \node at(-1.3,-1.6){箭头按 $Q_k>0$ 表示};
}
图\ref{fig:auxiliaryclausius}画出一份热交换的抵消关系。辅助机向 $T_k$ 送回的热量大小为 $Q_k$，它从参考热库取热；若 $Q_k<0$，相应方向全部反转。比较在全部循环结束后进行，辅助工质已恢复初态，无须再计其能量或熵的存量变化。

\subsection{实际边界温度与可逆替代路径}
熵差只要求初末为可定义的平衡态。实际传热时，系统内部可以存在温度梯度，此时不能随意选一个“平均温度”去除实际热量。熵收支中的温度必须对应所选边界处传递熵的温度；若把接触层一并纳入研究对象，可用外部热库温度计算熵输入，把接触层内的有限温差熵产生也计入 $S_{\rm gen}$。

对从恒温热库 $T_0$ 吸收总热 $Q$ 的封闭系统，
\[
 \Delta S_{\rm sys}-\frac Q{T_0}=S_{\rm gen}\ge0.
\]
左侧的系统熵差可沿任意可逆替代路径计算，热库熵差则按真实热量计算。替代路径只用于求熵差；计算实际热量、功和熵产生时，仍使用真实过程的边界条件。

若小块物体温度为 $T$，通过接触层从 $T_0>T$ 的热库吸收 $\delta Q$，物体在内部近似均温，则物体熵增为 $\delta Q/T$，热库熵变为 $-\delta Q/T_0$，两者之和
\[
 \delta S_{\rm gen}
 =\delta Q\frac{T_0-T}{TT_0}.
\]
温差趋于零时，每传递相同热量的熵产生趋于零；但传递速率通常也随温差降低。可逆极限与有限时间内的最大功率因而属于不同的过程要求。

\subsection{熵温图中的热量与路径方向}
对可逆过程 $\delta Q_{\rm rev}=T\dd S$。在以 $S$ 为横轴、$T$ 为纵轴的图上，沿曲线的有向面积给可逆热量。横向向右对应吸热，向左对应放热；等熵段为竖线，其热量为零。可逆循环满足 $\Delta U=0$，故封闭有向面积同时等于输出功。

\Diagram{fig:entropypatharea}{可逆路径下的有向面积与热量}{
 \draw[axis](0,0)--(5.5,0)node[right]{$S$};
 \draw[axis](0,0)--(0,3.4)node[above]{$T$};
 \fill[gray!15](1,0)--(1,1.3)--(4,2.5)--(4,0)--cycle;
 \draw[curve](1,1.3)--(4,2.5);
 \draw[->](2,1.7)--(3,2.1);
 \draw[guide](1,0)--(1,1.3);\draw[guide](4,0)--(4,2.5);
 \node[below]at(1,0){$S_A$};\node[below]at(4,0){$S_B$};
 \node at(2.5,.65){$\int_A^BT\dd S$};
}
图\ref{fig:entropypatharea}的面积只表示图中可逆路径的热量。不可逆绝热过程可以有 $\Delta S>0$，但实际 $Q=0$；把它的初末点连成一条线再求面积，算到的是另一个人为选择的可逆过程。相变等温段在熵温图上为水平线，面积为 $T\Delta S$，直接对应潜热。

